\chapter{The social state and its price}
\label{ch:socialstate}

\section{What was built}

Between 1961 and 1975 the Cuban state constructed the thing it is still known
for, and it did so fast enough that the international development literature
has been arguing about it ever since. The argument is worth having, but it
should not be allowed to obscure the outcomes, which are not in dispute and
which were achieved on an income per head that never rose above the middle of
the Latin American range.

\subsection{Health}

The starting point was bad and got worse before it got better: roughly half the
country's physicians emigrated between 1959 and 1963, leaving about three
thousand doctors for six and a half million people, most of them in Havana. The
response was structural rather than remedial. Medical education was expanded
past the point of surplus and kept there; a rural medical service was made
compulsory for new graduates; and the delivery model was rebuilt around the
\emph{policl\'inico}, a neighbourhood polyclinic responsible for a defined
population, with prevention, vaccination, maternal care and chronic-disease
follow-up as its main business rather than as an afterthought to curative
medicine. The family-doctor programme that put a physician and nurse on a city
block came later, in 1984, but it was the same logic pushed one level down.

The measurable results were large. Infant mortality, around 37 to 40 deaths per
thousand live births in the late 1950s by the incomplete registration of the
time, was under 20 by the mid-1970s and reached single digits in the
1980s.\unv{} Cuba eliminated polio in 1962, the first country in the Americas to
do so, through a mass campaign run by the block committees; measles,
diphtheria, tetanus and rubella followed. Life expectancy rose from around 62
years to over 73 by the mid-1980s, converging on the United States figure and
passing every country in the region except Costa Rica \citep{onei,feinsilver1993}.\unv{}

There are three honest qualifications, and none of them cancels the
achievement. Cuba's health indicators were already the best or nearly the best
in Latin America in 1958, so the revolution began from a high base and the
gains are partly a continuation of a pre-existing trend --- though the
convergence of rural with urban outcomes, which is the part the pre-existing
trend was not delivering, is entirely post-1959. Infant mortality is a
statistic sensitive to definitional practice and to a health system with strong
incentives to make it low; Cuba's aggressive antenatal management, including
very high rates of intervention in pregnancies flagged as risky, contributes to
the number in ways that a purely epidemiological reading misses. And the
achievement was bought with a share of national output that a country of Cuba's
income could only sustain because someone else was paying for the oil.

\subsection{Education}

Universal primary schooling was effectively achieved in the 1960s and universal
lower-secondary in the 1970s. University enrolment multiplied by roughly an
order of magnitude over three decades. Boarding secondary schools in the
countryside, the \emph{escuelas en el campo}, combined study with agricultural
work on a scale no other country attempted, moving hundreds of thousands of
adolescents into a state-run residential system for part of every year. Teacher
training was industrialised. Textbooks and materials were free.

The quality was real and independently confirmed later. When UNESCO's regional
laboratory tested third and fourth graders across thirteen Latin American
countries in 1997, Cuban children scored around 350 on a scale whose regional
mean was 250 and whose standard deviation was 50 --- roughly two standard
deviations above the region, in both language and mathematics, with the nearest
comparators, Brazil and Chile, near 260. The finding that matters more than the
average is distributional: pupils in Cuba's lowest-income schools outperformed
most upper-middle-class pupils in the rest of the region
\citep{carnoy2007}. That is the single most important number in this book's
argument for Part III, because it is the one that cannot be explained away by
subsidy: it says the country built a mass education system that compressed
rather than reproduced advantage, and that the system still worked after seven
years of the Special Period.

The cost side of education is that it was also, explicitly and by design, a
political formation system. The curriculum carried the ideology, access to
higher education was mediated by political record as well as academic
performance, and entire fields --- sociology, economics as practised elsewhere,
much of philosophy --- were closed or reduced to catechism for long periods.
Cuba produced excellent engineers, physicians and mathematicians, and almost no
economists, lawyers or social scientists trained to think about how a market or
a court works. Part III has to notice this: the country's human capital is
superb in exactly the areas a technical economy needs and thin in exactly the
areas an institutional transition needs.

\subsection{Women, and the Family Code}

Female labour-force participation rose from around 13 per cent in 1958 to over
30 per cent by the late 1970s and kept rising, driven by free childcare, free
education, and a Federation of Cuban Women that was simultaneously a mass
organisation, a service provider and a transmission belt for the party. The
Family Code of 1975 required, in law, that husbands share housework and
childcare equally --- a provision that was unenforceable and that everyone knew
was unenforceable, and that was nonetheless a declaration no other country in
the hemisphere had made. Abortion was legal, free and safe from 1965, decades
before most of Latin America. Cuban women's educational attainment overtook
men's and has stayed ahead.

What did not change was the composition of power. The Political Bureau remained
overwhelmingly male for decades, and the domestic division of labour changed
far less than the statute said it should.

\subsection{Race, declared solved}

The desegregation of 1959--61 was fast, real and enforced, and it is one of the
revolution's genuine moral achievements. What followed was a decision with long
consequences: in 1962 the leadership declared racial discrimination abolished
and the racial question closed, and independent Afro-Cuban organisations ---
including the societies and clubs that had been the black civil-society
tradition since the nineteenth century --- were dissolved along with everything
else. Race disappeared from the census categories of analysis and from
permissible public discussion for thirty years.

The effect was that a real gain was locked in and further progress stopped
being measurable. When race returned to Cuban public discussion in the 1990s it
did so because the dollar economy had made the disparity impossible to ignore:
remittances flowed overwhelmingly to white Cubans, because the 1960s emigration
had been overwhelmingly white, and the tourist sector hired by
\emph{buena presencia}. Cuba entered its market opening with a racial
distribution of foreign-currency access that it had spent thirty years
declining to measure \citep{delafuente2001}. Chapter~\ref{ch:premises} treats this as a hard
constraint rather than a footnote.

\section{What it cost}

The chapters of this book that describe Cuban social policy and the chapters
that describe Cuban repression are usually written by different people. They
describe the same state, in the same years, doing both things with the same
apparatus, and often through the same organisation --- the block committee that
ran the vaccination drive also kept the file on the neighbours.

\subsection{The camps}

Between 1965 and 1968 the Military Units to Aid Production, the UMAP, held
young men in agricultural labour camps in Camag\"uey under military discipline.
The categories swept up were religious believers --- Jehovah's Witnesses,
Seventh-Day Adventists, practising Catholics, including seminarians ---
homosexuals, and young men classed as antisocial, which in practice meant
long-haired, unemployed, or attached to the wrong music. Something like
twenty-five to thirty-five thousand people passed through them.\unv{} There
were deaths, beatings and suicides. The camps were closed in 1968 after
protests from Cuban intellectuals and international pressure, and the closure
was an admission, though the state has never made a full one; Castro said in
2010 that the persecution of homosexuals in those years was a ``great
injustice'' and that he bore responsibility for it.

\subsection{The five grey years}

The cultural opening of the early 1960s --- which was real, and which produced
the film institute, the Casa de las Am\'ericas, a serious literary magazine
culture and a poster art still worth looking at --- ended in 1971. The trigger
was the case of the poet Heberto Padilla, arrested in March 1971 after winning
a national prize for a book the writers' union had criticised, held for a
month, and released after a public self-criticism in which he denounced himself
and named his friends. The performance was so evidently coerced that it broke
the revolution's standing with the European and Latin American left; Sartre,
Vargas Llosa, Calvino, Paz, Sontag and dozens of others signed against it, and
several of them never came back.

The First National Congress on Education and Culture, a month later, made the
line official: art was a weapon of the revolution, homosexuals were unfit for
work with young people or for representing Cuba abroad, and religious belief
was incompatible with certain professions. What followed, 1971 to about 1976,
is known in Cuba as the \emph{quinquenio gris}, the five grey years, and its
mechanism was not mass arrest but \emph{parametraci\'on} --- the setting of
parameters. Writers, actors, academics and teachers were assessed against
ideological and personal criteria and, if they failed, removed from their
professions and reassigned to manual work. Careers were ended quietly, by
administrative letter, and many were never restored.

\subsection{Prison, exit, and the closing of the door}

The 1960s saw a rural insurgency in the Escambray mountains, fought from 1960
to 1965 and known officially as the struggle against bandits. It was defeated,
and the defeat included the forced relocation of thousands of peasant families
from the Escambray to new settlements in Pinar del R\'io, at the far end of the
island. The political prison population in the 1960s ran into the tens of
thousands at its peak; the \emph{plantados}, prisoners who refused the
re-education programme and its uniform, served terms of twenty and thirty
years.\unv{} No count has ever been published.

Emigration continued and its terms hardened. When Castro opened the port of
Camarioca in 1965, Cuban Americans came by boat; the resulting negotiation
produced the Freedom Flights, twice-daily charter flights from Varadero to
Miami that carried roughly 300,000 people between 1965 and 1973.\unv{} The
regime attached a price to using them. Applying to leave made a person a
\emph{gusano}, a worm; the applicant lost their job, was often assigned to
agricultural labour for the years-long wait, had their household inventoried
and confiscated on departure, and men of military age could not go at all.
Families were split by these rules for decades. The emotional economy of the
Cuban diaspora --- its intensity, its unwillingness to compromise, its
transmission across three generations --- is not fully explicable without
this policy, and Part III's proposals about the diaspora have to reckon with a
grievance that is specific and documented rather than diffuse.

\subsection{The Revolutionary Offensive of 1968}

In March 1968 the state nationalised what remained of small private business:
roughly 55,000 to 58,000 establishments, essentially all of them --- corner
shops, bars, food stalls, repair shops, laundries, street vendors,
carpenters.\unv{} The stated reason was that these were parasitic and a
breeding ground for the black market. The effect was to abolish the Cuban
service economy overnight, to hand its functions to a state distribution system
that had no capacity to perform them, and to eliminate in one law the entire
class of person who knew how to run a small enterprise.

This is the most consequential purely economic decision of the period, and it
is under-discussed relative to the agrarian reform because it did not involve
foreigners. Cuba spent the following fifty-eight years periodically
re-legalising, in narrow and reversible increments, categories of activity that
were legal and functioning in 1967: self-employment in 1978 and again in 1993,
private restaurants in 1993, private taxis, private rooms for rent,
agricultural markets in 1980 and again in 1994, and finally, in 2021, small and
medium private enterprises with legal personality. Every one of those steps was
presented as a bold reform. Every one of them was a partial restoration of
something the state had abolished on a single day in March 1968.

\begin{ledger}{the social state, 1961--1975}
\emph{Achieved.} Universal free health care with the rural--urban gap closed,
infant mortality more than halved, polio and measles eliminated, life
expectancy up more than ten years. Universal free education to secondary level
with measured attainment two standard deviations above the regional mean and
the strongest results at the bottom of the social distribution. Female
labour-force participation more than doubled, free childcare, legal abortion,
an equality statute ahead of the hemisphere. Public space desegregated by
enforcement rather than declaration. Social security extended to the whole
population.

\emph{Cost.} Forced-labour camps for believers, homosexuals and nonconformists,
1965--68. A cultural purge, 1971--76, conducted by administrative
reassignment rather than arrest and therefore leaving no record and no
redress. Political imprisonment on a scale never disclosed. Exit made
punishable, with confiscation and years of forced labour as the price of an
application, splitting families across three generations. All private
enterprise below the corner-shop level abolished in 1968, destroying the
country's service economy and its stock of entrepreneurial skill. Race declared
solved and then made unmeasurable for thirty years.
\end{ledger}
